% from ega3-7-fr.tex
\oldpage[III]{199}
\[
\xymatrix@C=18pt@R=32pt{
  {\begin{gathered}
    \bigl(\mathscr T_p(\mathscr O_Y)
      \otimes_{\mathscr O_Y}\mathscr M\bigr)|U\\
    =\mathscr T^U_p(\mathscr O_Y|U)
      \otimes_{\mathscr O_Y|U}(\mathscr M|U)
  \end{gathered}}
  \ar[r]\ar[d]
  &
  {\begin{gathered}
    \mathscr T^U_p(\mathscr M|U)\\
    =\bigl(\mathscr T_p(\mathscr M)\bigr)|U
  \end{gathered}}
  \ar[d]
  \\
  {\begin{gathered}
    \bigl(\mathscr T_p(\mathscr O_Y)
      \otimes_{\mathscr O_Y}\mathscr M\bigr)|V\\
    =\mathscr T^V_p(\mathscr O_Y|V)
      \otimes_{\mathscr O_Y|V}(\mathscr M|V)
  \end{gathered}}
  \ar[r]
  &
  {\begin{gathered}
    \mathscr T^V_p(\mathscr M|V)\\
    =\bigl(\mathscr T_p(\mathscr M)\bigr)|V
  \end{gathered}}
}
\]
est commutatif par (7.2.3.3).

Pour tout morphisme $g:Y'\to Y$ on a un homomorphisme canonique
\[
  \mathscr T_p(\mathscr O_Y)\otimes_{\mathscr O_Y}\mathscr O_{Y'}
  \longrightarrow
  \mathscr T^{Y'}_p(\mathscr O_{Y'}),
  \tag{7.7.3.2}\label{III.7.7.3.2.fr}
\]
qui n'est autre que le cas particulier de (6.7.11.2) (pour les aboutissements)
dans le cas où $S=Y$, $v_1=f$, $v_2=1_Y$, $\mathscr P^{(2)}_\bullet$ réduit au
seul terme $\mathscr M$ de degré 0.

Lorsque $Y=\operatorname{Spec}(A)$, $Y'=\operatorname{Spec}(A')$ sont affines,
(7.7.3.2) n'est autre que l'homomorphisme de faisceaux correspondant à
l'homomorphisme canonique de $A'$-modules défini dans (7.2.2)
\[
  T^A_p(A)\otimes_A A'
  \longrightarrow
  T^{A'}_p(A')=T^A_p(A')
\]
comme il résulte aisément de (6.7.11) (car dans le cas envisagé, on peut prendre
$\mathfrak U'^{(i)}=u_i^{-1}(\mathfrak U^{(i)})$ dans (6.7.11)).
\end{env}

\begin{env}[7.7.4]
Lorsque $f$ est un morphisme \emph{propre}, $Y=\operatorname{Spec}(A)$ un schéma
affine noethérien et $\mathscr P_\bullet$ un complexe de
$\mathscr O_X$-Modules cohérents et $Y$-plats \emph{limité inférieurement}, on a
vu (6.10.5) que l'on peut écrire à un isomorphisme près,
\[
  \mathscr T_p(\mathscr M)
  =\mathscr H_p
    \bigl(\mathscr L_\bullet\otimes_{\mathscr O_Y}\mathscr M\bigr),
  \qquad
  \mathscr L_\bullet=\widetilde{L_\bullet},
\]
où $L_\bullet$ est un complexe de $A$-modules \emph{libres de type fini limité
inférieurement}; le foncteur $\mathscr T_\bullet$ est donc du type qui a été
étudié en détail dans (7.4) et (7.6). Nous allons traduire les résultats de cette
étude :
\end{env}

\begin{theorem}[7.7.5]
Soient $Y$ un préschéma localement noethérien, $(U_\alpha)$ un recouvrement de
$Y$ formé d'ouverts affines, $f:X\to Y$ un morphisme propre,
$\mathscr P_\bullet$ un complexe de $\mathscr O_X$-Modules cohérents et $Y$-plats
limité inférieurement. Le foncteur homologique $\mathscr T_\bullet(\mathscr M)$
défini par (7.7.1.1) possède alors les propriétés suivantes :

\medskip
\noindent
\textnormal{I) (La propriété de semi-continuité).}\footnote{Un cas particulier de
ce théorème se trouve déjà dans la note [3] de Chow-Igusa. La propriété de
semi-continuité a été découverte, dans le cadre des espaces analytiques (et sous
des hypothèses assez particulières), par Kodaira-Spencer (« On the variations of
almost-complex structures », \emph{Algebraic Geometry and Topology, A Symposium in
honor of S. Lefschetz}, Princeton Series n\textsuperscript{o} 12, p. 139--150,
Princeton, 1957) et la version générale démontrée par Grauert [5].}
\emph{La fonction}
\[
  y\longmapsto
  d_p(y)
  =
  \operatorname{rang}_{\kappa(y)}
    \mathscr T_p^{\kappa(y)}(\kappa(y))
  \tag{7.7.5.1}\label{III.7.7.5.1.fr}
\]
\emph{est semi-continue supérieurement.}

